\documentclass[11pt,a4paper]{article}

\usepackage[utf8]{inputenc}
\usepackage[T1]{fontenc}
\usepackage{lmodern}
\usepackage{geometry}
\geometry{a4paper, margin=1in}

\usepackage{titlesec}
\usepackage{hyperref}
\usepackage{graphicx}
\usepackage{xcolor}
\usepackage{listings}
\usepackage{enumitem}

% Setup links
\hypersetup{
    colorlinks=true,
    linkcolor=blue,
    filecolor=magenta,      
    urlcolor=cyan,
    pdftitle={The Complexity of Credit Risk Data Engineering},
}

% Setup code blocks
\definecolor{codegray}{rgb}{0.5,0.5,0.5}
\definecolor{codepurple}{rgb}{0.58,0,0.82}
\definecolor{backcolour}{rgb}{0.95,0.95,0.92}

\lstdefinestyle{mystyle}{
    backgroundcolor=\color{backcolour},   
    commentstyle=\color{codegray},
    keywordstyle=\color{blue}\bfseries,
    numberstyle=\tiny\color{codegray},
    stringstyle=\color{codepurple},
    basicstyle=\ttfamily\footnotesize,
    breakatwhitespace=false,         
    breaklines=true,                 
    captionpos=b,                    
    keepspaces=true,                 
    numbers=none,                    
    showspaces=false,                
    showstringspaces=false,
    showtabs=false,                  
    tabsize=2
}

\lstset{style=mystyle}

\title{\Huge \textbf{The Complexity of Credit Risk Data Engineering}}
\author{Context: Home Credit Dataset $\cdot$ Kaggle 1st Place Insights}
\date{\today}

\begin{document}

\maketitle

\section{The Real Challenge: Heterogeneous, Time-Traveling Data}

As the 1st Place Kaggle Winner noted, credit underwriting is one of the most complex domains for machine learning because the data is:
\begin{itemize}
    \item \textbf{Heterogeneous:} It mixes static demographics, changing incomes, point-in-time bureau snapshots, and high-frequency transaction logs.
    \item \textbf{Collected over different time frames:} A borrower's income might have been recorded 2 years ago, their bureau score 1 month ago, and their last late payment 3 days ago.
    \item \textbf{Subject to changing collection environments:} A bank might switch from asking ``Monthly Income'' to ``Annual Income'' in the middle of a database's life, silently breaking standard logic.
\end{itemize}

\subsection{The Target Variable Trap}
Defining the target (\texttt{1} = Default, \texttt{0} = Paid) is arguably the hardest part of credit ML. 
\begin{itemize}
    \item Do you define a default as 30 days past due (DPD), 60 DPD, or 90 DPD? 
    \item If someone is 30 days late but then pays it all back (a ``cure''), are they a \texttt{1} or a \texttt{0}?
    \item What is the performance window? Do we wait 12 months to label a loan, or 24 months?
\end{itemize}

The Home Credit dataset is celebrated because they solved these domain complexities \textit{before} releasing the data, making it ``leak-free'' (meaning you cannot accidentally use future data to predict the past).

\section{Smart Features vs. Base Algorithms}

The winning Kaggle philosophy relies on two pillars:
\begin{enumerate}
    \item \textbf{Good set of smart features} (Domain Knowledge)
    \item \textbf{Diverse base algorithms} (Ensembling LightGBM, XGBoost, CatBoost, Neural Nets)
\end{enumerate}

While anyone can run a LightGBM model, the magic is in the feature engineering. The winner used \textbf{$\sim$700 engineered features}, primarily generated by aggregating many-to-one tables.

\subsection{The ``Many-to-One'' Aggregation Strategy}
For every single borrower in \texttt{application\_train.csv} (1 row per user), there are dozens of rows in other tables:
\begin{itemize}
    \item 5 previous applications in \texttt{previous\_application.csv}
    \item 30 historical credit accounts in \texttt{bureau.csv}
    \item 120 past payments in \texttt{installments\_payments.csv}
\end{itemize}

\textbf{You cannot feed 120 rows into a standard ML model to predict 1 outcome.} You must flatten them into ``smart aggregates.''

\subsubsection{Examples of Smart Aggregates (The 700 Features):}
\begin{itemize}
    \item \textbf{Averages:} What is their \textit{average} past due amount across all previous loans?
    \item \textbf{Max/Mins:} What is the \textit{maximum} number of days they were ever late?
    \item \textbf{Recency Weighted:} How many times were they late \textit{in the last 3 months} vs. \textit{in the last 2 years}? (Recent behavior is always more predictive).
    \item \textbf{Ratios:} (Total past payments made) $\div$ (Total past payments required).
\end{itemize}

\section{Data Cleaning: The Unseen Foundation}

Before you can aggregate, you have to clean. Credit data is notorious for silent errors. 

\subsection{The Sentinel Value Problem}
In our notebook, we saw:
\begin{lstlisting}[language=Python]
frame["DAYS_EMPLOYED_ANOM"] = (frame["DAYS_EMPLOYED"] == 365243).astype(int)
frame.loc[frame["DAYS_EMPLOYED"] == 365243, "DAYS_EMPLOYED"] = np.nan
\end{lstlisting}
\begin{itemize}
    \item A value of \texttt{365243} days employed equals \textbf{1,000 years}. 
    \item This was not an immortal applicant; it was Home Credit's system using \texttt{365243} as a ``sentinel value'' to mean ``Unemployed.''
    \item If you fed this directly to a model, the algorithm would think this group is extremely employed and stable, completely reversing the actual risk profile (unemployed people are highly risky).
\end{itemize}

\subsection{Safe Math in Finance}
\begin{lstlisting}[language=Python]
def safe_divide(a, b):
    if isinstance(b, pd.Series):
        b = b.replace({0: np.nan})
    return a / b
\end{lstlisting}
In credit data, dividing by zero happens constantly (e.g., trying to find \texttt{Income $\div$ Debt} when they have \$0 debt). A standard Python \texttt{ZeroDivisionError} will crash your entire pipeline. \texttt{safe\_divide} forces these to \texttt{NaN} (Missing), which Gradient Boosting models handle natively.

\section{The Goto Financial / Fintech Context}

For a modern fintech dealing with thin-file borrowers, the aggregation strategy shifts from formal banking tables to behavioral logs:

\begin{enumerate}
    \item \textbf{Velocity Aggregates:} Instead of \texttt{bureau.csv}, you aggregate GoPay logs. \textit{How many times did this user transfer money out immediately after receiving a transfer in?} (High velocity = bridging loans/distress).
    \item \textbf{Time-of-Day Features:} \textit{What percentage of their transactions happen between 12 AM and 4 AM?} (Late night high-frequency transactions often correlate with risky behavior like online gambling).
    \item \textbf{App Telemetry:} \textit{How long did they spend reading the terms and conditions? Did they copy-paste their ID number or type it out?} (Typing it from memory often indicates fraud rings).
\end{enumerate}

The 1st place Kaggle lesson remains true: the algorithm is a commodity. The competitive moat is how creatively you can aggregate thousands of messy historical data points into a single, clean user profile.

\end{document}
